% Ported from sections/03_discussion.md (v5, APPROVED 2026-09-18).
% Final-manuscript Sec.~IV.
\section{Discussion}
\label{sec:discussion}

\subsection{Why the island is the ground state}
\label{sec:island-ground}

The search finds an island ground state in both models, with the flat film sitting 0.19~\eVA{} higher
for Fe/MgO and 0.15~\eVA{} higher for Fe-B/MgO. The electronic-structure analysis (Supplementary
Material) indicates why. In the flat monolayer every metal atom is registry-locked directly atop an
oxygen of the MgO surface---the registry of the reference construction, and the one determined
experimentally for the first monolayer of Fe on MgO(001) by LEED I--V analysis \cite{urano1988} and
used in first-principles models of the Fe$|$MgO$|$Fe interface \cite{butler2001}---maximising Fe--O
orbital overlap and pushing the Fe $d$-band centre down to $-0.23$~eV. The island abandons that
registry: the number of metal atoms registered directly atop an oxygen drops from 25 (of 25) to 9 (of
25), the island's interface Fe $d$-band centre rises to $+0.51$~eV ($+0.60$~eV averaged over all
island Fe), and the film becomes electronically ``quieter'' (lower DOS at the Fermi level, weaker
spin polarisation). The island's energy gain is therefore dominated by \textbf{reduced forced
interfacial coupling and restored metal cohesion}: the registry-locked monolayer spends its bonding
on Fe--O contacts while forgoing the three-dimensional Fe--Fe coordination available to a cluster,
and the island reverses that trade.

It is \textbf{not} a lattice-strain effect. The in-plane mismatch of this interface (3.6\,\% of the
MgO lattice against the Fe lattice, Sec.~\ref{sec:models}) is carried by the \textbf{substrate},
which is built in that compressed state and held fixed, not by the film; the Fe film sits at its own
equilibrium lattice constant and is not strained in-plane, so there is no film strain for the island
to relieve. The mechanism is consistent with the weak Fe--MgO coupling found in first-principles
treatments of the interface \cite{butler2001}, and with the experimental observation that a flat
monolayer requires low-temperature or slow deposition while room-temperature growth gives 3D islands
at the same one-monolayer coverage \cite{fahsold2000,torelli2009,reitinger2007}.

\subsection{Why boron stabilises the flat film}
\label{sec:boron}

Boron lowers the flat-state energy by 0.040~\eVA, bringing the flat configuration closer to the
ground state. Two observations frame the mechanism. First, boron does \textbf{not} bond to the MgO
interface---it remains inside the metal film (a single B--O contact in one of 72 windowed Fe-B
structures)---so its effect is not interfacial. Second, the shift is not carried by one favourable
search: the Fe-B model's six per-search flat minima lie between 0.1493 and 0.2481~\eVA, five of them
below the Fe host's median of 0.2185~\eVA, and 94\,\% of resampled replicates agree in sign
(Sec.~\ref{sec:flat}). Both observations point to an intrinsic, film-internal role for boron rather
than an interfacial or statistical one.

\textbf{Confusion principle and amorphous formation.} The flat film can be read as the disordered,
amorphous-like configuration and the island as the ordered, crystalline-like one. Under this reading,
the stabilisation of the flat state by added elements follows the \textbf{confusion principle} of
metallic-glass formation \cite{greer1993}: the more elements in an alloy, the harder it is for the
alloy to select a viable crystal structure, and the greater the tendency toward glass (amorphous)
formation.

Our results are partly consistent with this. Adding boron lowers the flat-basin energy
($0.1888 \to 0.1493$~\eVA), i.e.\ the added element stabilises the disordered configuration relative
to the ordered one, in the direction the principle predicts. The present models test this for a
single added metalloid and cannot decide between the competing readings of the principle---whether it
acts through the \emph{presence of an added species} by a film-internal mechanism, or through the
\emph{number} of species, which would require the same host with two different additions to separate.
We therefore read our result as consistent with the confusion principle acting through the added
metalloid and leave the element-count reading untested here (see Sec.~\ref{sec:limitations}).

We note, however, that the boron-containing film retains the island as its ground state---the
confusion principle stabilises the flat state but does not, in this small model system, fully
suppress the ordered configuration. The precise origin of boron's effect---whether it lowers the flat
basin's energy or raises the island's---is not resolved by the present data and is a natural target
for the re-relaxation and further analysis.

\subsection{Implications for MTJ stacks}
\label{sec:mtj}

The performance of MgO-based magnetic tunnel junctions rests on the structural quality of the
metal/oxide interface. The giant tunnel magnetoresistance of single-crystal Fe/MgO/Fe junctions
arises from coherent spin-polarised tunnelling across a lattice-matched interface, with the residual
mismatch accommodated by interfacial dislocations and the growth conditions chosen to minimise them
\cite{yuasa2004}. The present results connect to this in two ways. First, they show that a perfectly
flat, registry-locked metal layer maximises Fe--O bonding at the expense of metal coordination and is
therefore energetically penalised relative to a clustered film---a consideration for interface
engineering in MgO-based tunnel junctions, whose barrier must stay flat and coherently matched for
the tunnelling magnetoresistance to reach its high values \cite{yuasa2004}. Second, they show that
boron acts to stabilise the flat configuration without bonding to the interface, consistent with the
picture of boron as a film-internal agent that promotes a flat, well-wetting interface.
These are trend-level, model-system conclusions; quantitative transfer to a device stack would
require the converged relaxations and a fuller treatment of the interface.

\subsection{Limitations}
\label{sec:limitations}

One scope bound follows from the design rather than from the methods. The paper compares a single
host with and without one added element, boron. Nothing here separates the effect of \emph{boron}
from the effect of \emph{an added element as such}; that would need a second chemically different
addition in the same host, which is not analysed (the element-count reading of the confusion
principle is therefore untested---Sec.~\ref{sec:boron}). What the design does exclude is a trivial
explanation from the construction: the two models share the lattice constant, substrate,
reference-layer geometry and candidate-generation schedule, and differ only by the boron
(Sec.~\ref{sec:scope}).

The results are qualitative/trend-level: the structures are not DFT-converged minima (residual
forces $\sim$1--2~eV/\AA), the models are single-layer slabs at $\Gamma$-point sampling, and the
flat/island split uses a chosen $\dZ$ threshold. The two models also have unequal seed counts (13
completed searches for Fe/MgO against 6 for Fe-B/MgO) and their atom counts differ (75 and 78), so
sampling fractions are compared only as trends, not as like-for-like populations. The flat state is
described as a higher-energy configuration, not a proven metastable state. A full re-relaxation of
representative structures is necessary to place these conclusions on converged minima.

Two features of the model itself bound the interpretation. First, the \textbf{strain convention is
the inverse of the experimental stack}: here the in-plane lattice is the Fe lattice constant and the
\textbf{MgO substrate} is the compressed component, held fixed, whereas in a real junction the bulk
MgO imposes its lattice on a thin Fe film, which absorbs the mismatch as in-plane strain and relieves
it through interfacial dislocations \cite{yuasa2004}. The film in this work is therefore unstrained,
and this model does not represent the strained-film situation. Second, both phases are built on a
body-centred-cubic Fe lattice, while the experimentally reported structure of ultrathin Fe on
MgO(001) is body-centred tetragonal below about 10~\AA{} \cite{urano1988}; the island branch
($\dZ \approx 1$--$5.6$~\AA) lies in that regime (Sec.~\ref{sec:scope}). The flat--island comparison
is therefore a trend obtained within one fixed lattice model, and extending it to thicker,
experimentally strained films would require a different construction.

\textbf{The search-level statistic rests on 19 searches, and it is asymmetric.} The comparison pools
13 completed searches of the Fe/MgO model against 6 of the Fe-B/MgO model. The permutation test
enumerates all 27\,132 partitions of that pool, so its resolution floor is $p = 3.7\times10^{-5}$ and
the reported $p = 0.0444$ is not limited by the number of searches; it would, however, tighten with
more completed Fe-B/MgO searches, since the smaller group sets the width of the null. Two further
bounds follow from the same asymmetry. The boron-free interface statement (Sec.~\ref{sec:flat}) is
measured on the 72 Fe-B/MgO structures inside the low-energy window, which come from \textbf{four} of
the six searches---so its effective number of independent observations is closer to four than to
72---and the Fe-B/MgO motif counts (Sec.~\ref{sec:landscape}) rest on two (flat) and three (island)
searches against five and four for Fe/MgO. Statements resting on those sets are correspondingly
weaker evidence than the flat-basin energies of Table~\ref{tab:flatbasin}, which use every search of
both models.
